\section{Finite-conductor energy escape and a layerwise boundary}
\label{sec:finite-conductor-barriers}

The almost-all-shift theorem uses genuinely nonlocal information about prime
correlations.  We next quantify why simply fitting the prime fluctuation by a
periodic observable does not recover comparable $L^2$ energy.

\subsection{Orthogonal projection onto periodic observables}
\label{subsec:periodic-projection}

Let
\begin{equation}\label{eq:dyadic-integer-interval}
 I_N=\{n\in\Z:N<n\le2N\}
\end{equation}
and equip $\ell^2(I_N)$ with counting measure.  Let $V_q(N)$ be the subspace
of restrictions to $I_N$ of $q$-periodic complex functions, and let $P_q$ be
the orthogonal projection onto $V_q(N)$.

\begin{theorem}[Finite-conductor energy escape]
\label{thm:finite-conductor-energy-escape}
Fix $0<\delta<1$.  Uniformly for integers
\begin{equation}\label{eq:energy-q-range}
 1\le q\le N^{1-\delta},
\end{equation}
one has
\begin{equation}\label{eq:periodic-projection-bound}
 \|P_q(\Lambda-1)\|_{\ell^2(I_N)}^2
 \ll_\delta N\frac q{\varphi(q)}
 \ll_\delta N\log\log(3q).
\end{equation}
On the other hand,
\begin{equation}\label{eq:prime-fluctuation-energy}
 \|\Lambda-1\|_{\ell^2(I_N)}^2\sim N\log N.
\end{equation}
Consequently, uniformly in \eqref{eq:energy-q-range},
\begin{align}
 \frac{\|P_q(\Lambda-1)\|_2^2}{\|\Lambda-1\|_2^2}
 &\ll_\delta\frac{\log\log(3q)}{\log N}=o_\delta(1),
 \label{eq:captured-energy-ratio}\\
 \inf_{b\in V_q(N)}\|\Lambda-1-b\|_2^2
 &=(1-o_\delta(1))N\log N.
 \label{eq:periodic-best-approximation}
\end{align}
\end{theorem}

\begin{proof}
For $a\pmod q$, let
\begin{equation}\label{eq:residue-class-count}
 N_a=\#\{n\in I_N:n\equiv a\pmod q\}.
\end{equation}
The range \eqref{eq:energy-q-range} gives
$N_a=N/q+O(1)\asymp_\delta N/q$ uniformly for sufficiently large $N$.
For any $f\in\ell^2(I_N)$, the projection is the residue-class average, so
\begin{equation}\label{eq:projection-exact-formula}
 \|P_qf\|_2^2
 =\sum_{a\, (\mathrm{mod}\,q)}
  \frac1{N_a}
  \left|\sum_{\substack{n\in I_N\\n\equiv a\, (\mathrm{mod}\,q)}}
  f(n)\right|^2.
\end{equation}

Let $\Lambda_{\mathrm{pr}}(n)=\log n$ when $n$ is prime and zero otherwise.
If $(a,q)=1$, the Brun--Titchmarsh inequality and
$\log(N/q)\ge\delta\log N$ imply
\begin{equation}\label{eq:BT-weighted-row}
 \sum_{\substack{N<p\le2N\\p\equiv a\, (\mathrm{mod}\,q)}}\log p
 \ll_\delta\frac N{\varphi(q)}.
\end{equation}
We use here only the classical Brun--Titchmarsh inequality; the other
standard prime-number estimates in this proof may be found, for example, in
\citet{MontgomeryVaughan2007}.
Thus the residue sum of $\Lambda_{\mathrm{pr}}-1$ has absolute value
$O_\delta(N/\varphi(q))$.  If $(a,q)>1$, no prime $p>N>q$ lies in that
class, and the residue sum is simply $-N_a$.  Substitution into
\eqref{eq:projection-exact-formula} yields
\begin{equation}\label{eq:prime-only-projection}
 \|P_q(\Lambda_{\mathrm{pr}}-1)\|_2^2
 \ll_\delta N\frac q{\varphi(q)}+N
 \ll_\delta N\frac q{\varphi(q)}.
\end{equation}

The prime-power remainder
$R_{\mathrm{pp}}=\Lambda-\Lambda_{\mathrm{pr}}$ satisfies
\begin{equation}\label{eq:prime-power-L2}
 \|R_{\mathrm{pp}}\|_2^2
 \le\sum_{\substack{N<p^k\le2N\\k\ge2}}(\log p)^2
 \ll N^{1/2}(\log N)^2.
\end{equation}
Orthogonal projection is a contraction.  Combining
\eqref{eq:prime-only-projection} and \eqref{eq:prime-power-L2}, using
$\|u+v\|^2\le2\|u\|^2+2\|v\|^2$, proves the first bound in
\eqref{eq:periodic-projection-bound}.  The standard estimate
$q/\varphi(q)\ll\log\log(3q)$ proves the second.

The prime number theorem and partial summation give
\[
 \sum_{n\in I_N}\Lambda(n)\sim N,
 \qquad
 \sum_{n\in I_N}\Lambda(n)^2\sim N\log N,
\]
which proves \eqref{eq:prime-fluctuation-energy}.  Finally, the orthogonal
decomposition
\[
 \inf_{b\in V_q(N)}\|\Lambda-1-b\|_2^2
 =\|(I-P_q)(\Lambda-1)\|_2^2
 =\|\Lambda-1\|_2^2-\|P_q(\Lambda-1)\|_2^2
\]
gives \eqref{eq:captured-energy-ratio} and
\eqref{eq:periodic-best-approximation}.
\end{proof}

\begin{remark}[Exact boundary of the theorem]
The space $V_q(N)$ is a linear finite-conductor model.  A small projection
onto $V_q(N)$ says nothing by itself about bilinear forms, adaptive weights,
growing collections of moduli, or nonlinear sieve information.  Thus
\cref{thm:finite-conductor-energy-escape} is not a general impossibility
theorem for analytic number theory.
\end{remark}

\subsection{The outer divisor slice already sees shifted-prime M\"obius}
\label{subsec:outer-divisor-slice}

The exact defect can also be written
\begin{equation}\label{eq:defect-hyperbolic-ledger}
 \cT_{h,>D}(X)
 =-\sum_{D<d\le2X}\mu(d)\log d
   \sum_{X<dm\le2X}\Lambda(dm+h).
\end{equation}
If $d>X$, the inner variable is forced to be $m=1$.  Therefore the outermost
slice is exactly
\begin{equation}\label{eq:outermost-defect-slice}
 \cO_h(X)
 =-\sum_{X<d\le2X}\mu(d)\log d\,\Lambda(d+h).
\end{equation}
For fixed $h$, removing target prime powers of exponent at least two gives
\begin{equation}\label{eq:outer-slice-shifted-primes}
 \cO_h(X)
 =-\sum_{X+h<p\le2X+h}
   \mu(p-h)\log(p-h)\log p
 +O_h\bigl(X^{1/2}(\log X)^2\bigr).
\end{equation}

The neighboring unweighted assertion
\begin{equation}\label{eq:shifted-prime-mobius-folklore}
 \sum_{p\le X}\mu(p+h)=o(\pi(X))
\end{equation}
for a fixed nonzero shift is a folklore open problem.  Averaged-shift forms
of \eqref{eq:shifted-prime-mobius-folklore} are known
\citep{Lichtman2022,LichtmanTeravainen2022}.  Formula
\eqref{eq:outer-slice-shifted-primes} explains why a proposed proof that
bounds every divisor slice separately already meets a difficult
fixed-shift correlation at the outer edge.

\begin{remark}[No inference for an individual layer]
The smallness of the complete sum in
\cref{thm:complete-defect-almost-all} does not imply that
\eqref{eq:outermost-defect-slice} is small.  Other ranges of $d$ in
\eqref{eq:defect-hyperbolic-ledger} may cancel it.  The observation above is
a warning against layerwise absolute-value arguments, not a deduction from
the complete-tail theorem.
\end{remark}
